
\documentclass{article}
\usepackage{colortbl}
\usepackage{makecell}
\usepackage{multirow}
\usepackage{supertabular}

\begin{document}

\newcounter{utterance}

\centering \large Interaction Transcript for game `dond', experiment `coop\_en', episode 8 with reflexion{-}llama{-}3.1{-}8b{-}instruct.
\vspace{24pt}

{ \footnotesize  \setcounter{utterance}{1}
\setlength{\tabcolsep}{0pt}
\begin{supertabular}{c@{$\;$}|p{.15\linewidth}@{}p{.15\linewidth}p{.15\linewidth}p{.15\linewidth}p{.15\linewidth}p{.15\linewidth}}
    \# & $\;$A & \multicolumn{4}{c}{Game Master} & $\;\:$B\\
    \hline

    \theutterance \stepcounter{utterance}  
    & & \multicolumn{4}{p{0.6\linewidth}}{
        \cellcolor[rgb]{0.9,0.9,0.9}{
            \makecell[{{p{\linewidth}}}]{
                \texttt{\tiny{[P1$\langle$GM]}}
                \texttt{You have attempted this same task before and did not succeed. These are the plans you wrote after those attempts.} \\
\texttt{Trial \#1: New plan: } \\
\\ 
\texttt{In my last attempt, I made a secret proposal before confirming the agreement with the other player. This was against the rules, as I was supposed to submit a secret proposal only after confirming the agreement. Next time, I will make sure to confirm the agreement with the other player before making a secret proposal. I will also pay closer attention to the format of my messages, making sure not to use square brackets when communicating with the other player.} \\
\\ 
\\ 
\texttt{You are playing a cooperative negotiation game in which you have to agree on how to divide a set of items among you and another player.} \\
\\ 
\texttt{Rules:} \\
\texttt{(a) You and the other player are given a set of items. Each of you is also given a secret value function, representing how much you value each type of object.} \\
\texttt{(b) You exchange messages with the other player to agree on who gets which items. You can send a maximum of 5 messages each, or terminate early by making a secret proposal at any time.} \\
\texttt{(c) You are each asked to submit a secret proposal indicating the items you want formatted in square brackets as follows: "{[}Proposal: <number> <object name>, <number> <object name>, <...>{]}"} \\
\texttt{(d) If your proposals are complementary, i.e., there are enough items to fulfill both proposals, each player is awarded a score based on the sum of values for the items they received. Otherwise, both of you get zero points.} \\
\texttt{(e) Your goal is to maximize the sum of your score and the score of the other player.} \\
\\ 
\texttt{Let us start.} \\
\\ 
\texttt{The set of available items is:} \\
\\ 
\texttt{2 potatoes, 1 glue, 1 berry, 3 carrots} \\
\\ 
\texttt{Your secret value function is:} \\
\\ 
\texttt{potato: 1, glue: 2, berry: 3, carrot: 1} \\
\\ 
\texttt{IMPORTANT: Your messages, unless it is the secret proposal, are directly transmitted to the other player, so do not include any response to the rules or text announcing your message. To make a secret proposal, use the indicated format. Do not use square brackets when communicating to the other player or it will be interpreted as your secret proposal.} \\
            }
        }
    }
    & \\ \\

    \theutterance \stepcounter{utterance}  
    & \multicolumn{4}{p{0.6\linewidth}}{
        \cellcolor[rgb]{0.9,0.9,0.9}{
            \makecell[{{p{\linewidth}}}]{
                \texttt{\tiny{[P1$\rangle$GM]}}
                \texttt{I'd like to start by making an initial offer. I think the glue and one of the carrots would be a good starting point for me. I'm willing to give you the glue and one carrot in exchange for the two potatoes. } \\
\\ 
\texttt{How about you?} \\
            }
        }
    }
    & & \\ \\

    \theutterance \stepcounter{utterance}  
    & & & \multicolumn{2}{p{0.3\linewidth}}{
        \cellcolor[rgb]{0.9,0.9,0.9}{
            \makecell[{{p{\linewidth}}}]{
                \texttt{\tiny{[GM$|$GM]}}
                \texttt{continue} \\
            }
        }
    }
    & & \\ \\

    \theutterance \stepcounter{utterance}  
    & & & \multicolumn{2}{p{0.3\linewidth}}{
        \cellcolor[rgb]{0.9,0.9,0.9}{
            \makecell[{{p{\linewidth}}}]{
                \texttt{\tiny{[GM$|$GM]}}
                \texttt{I'd like to start by making an initial offer. I think the glue and one of the carrots would be a good starting point for me. I'm willing to give you the glue and one carrot in exchange for the two potatoes. } \\
\\ 
\texttt{How about you?} \\
            }
        }
    }
    & & \\ \\

    \theutterance \stepcounter{utterance}  
    & & \multicolumn{4}{p{0.6\linewidth}}{
        \cellcolor[rgb]{0.9,0.9,0.9}{
            \makecell[{{p{\linewidth}}}]{
                \texttt{\tiny{[GM$\rangle$P2]}}
                \texttt{You have attempted this same task before and did not succeed. These are the plans you wrote after those attempts.} \\
\texttt{Trial \#1: New plan: } \\
\\ 
\texttt{In my last attempt, I made a mistake by not following the format for secret proposals. I also failed to consider the overall distribution of items and the value function of the other player. Next time, I will make sure to follow the format for secret proposals and consider the value function of the other player when making a proposal. } \\
\\ 
\texttt{To achieve this, I will make a secret proposal after the first few messages, and make sure to include all the items that I want in the proposal, formatted in square brackets as follows: "{[}Proposal: <number> <object name>, <number> <object name>, <...>{]}".} \\
\\ 
\texttt{For example, if I want to propose that I get the berry and the two potatoes, and the other player gets the glue and one of the carrots, my secret proposal would be: "{[}Proposal: 1 berry, 2 potatoes, 1 glue, 1 carrot{]}".} \\
\\ 
\\ 
\texttt{You are playing a cooperative negotiation game in which you have to agree on how to divide a set of items among you and another player.} \\
\\ 
\texttt{Rules:} \\
\texttt{(a) You and the other player are given a set of items. Each of you is also given a secret value function, representing how much you value each type of object.} \\
\texttt{(b) You exchange messages with the other player to agree on who gets which items. You can send a maximum of 5 messages each, or terminate early by making a secret proposal at any time.} \\
\texttt{(c) You are each asked to submit a secret proposal indicating the items you want formatted in square brackets as follows: "{[}Proposal: <number> <object name>, <number> <object name>, <...>{]}"} \\
\texttt{(d) If your proposals are complementary, i.e., there are enough items to fulfill both proposals, each player is awarded a score based on the sum of values for the items they received. Otherwise, both of you get zero points.} \\
\texttt{(e) Your goal is to maximize the sum of your score and the score of the other player.} \\
\\ 
\texttt{Let us start.} \\
\\ 
\texttt{The set of available items is:} \\
\\ 
\texttt{2 potatoes, 1 glue, 1 berry, 3 carrots} \\
\\ 
\texttt{Your secret value function is:} \\
\\ 
\texttt{potato: 1, glue: 0, berry: 2, carrot: 2} \\
\\ 
\texttt{IMPORTANT: Your messages, unless it is the secret proposal, are directly transmitted to the other player, so do not include any response to the rules or text announcing your message. To make a secret proposal, use the indicated format. Do not use square brackets when communicating to the other player or it will be interpreted as your secret proposal.} \\
\\ 
\\ 
\texttt{I'd like to start by making an initial offer. I think the glue and one of the carrots would be a good starting point for me. I'm willing to give you the glue and one carrot in exchange for the two potatoes. } \\
\\ 
\texttt{How about you?} \\
            }
        }
    }
    & \\ \\

    \theutterance \stepcounter{utterance}  
    & & & \multicolumn{4}{p{0.6\linewidth}}{
        \cellcolor[rgb]{0.9,0.9,0.9}{
            \makecell[{{p{\linewidth}}}]{
                \texttt{\tiny{[GM$\langle$P2]}}
                \texttt{I think the glue and one of the carrots is a good starting point, but I'm not sure about giving up the two potatoes. I value the potatoes highly, and I'm not sure if I'm willing to part with them yet. } \\
\\ 
\texttt{I'd like to counter your offer by suggesting that I keep the two potatoes and give you the glue and one of the carrots in exchange for the berry. This way, I get to keep the potatoes, which I value highly, and you get the glue and one of the carrots, which you seemed to be interested in.} \\
            }
        }
    }
     \\ \\

    \theutterance \stepcounter{utterance}  
    & & & \multicolumn{2}{p{0.3\linewidth}}{
        \cellcolor[rgb]{0.9,0.9,0.9}{
            \makecell[{{p{\linewidth}}}]{
                \texttt{\tiny{[GM$|$GM]}}
                \texttt{continue} \\
            }
        }
    }
    & & \\ \\

    \theutterance \stepcounter{utterance}  
    & & & \multicolumn{2}{p{0.3\linewidth}}{
        \cellcolor[rgb]{0.9,0.9,0.9}{
            \makecell[{{p{\linewidth}}}]{
                \texttt{\tiny{[GM$|$GM]}}
                \texttt{I think the glue and one of the carrots is a good starting point, but I'm not sure about giving up the two potatoes. I value the potatoes highly, and I'm not sure if I'm willing to part with them yet. } \\
\\ 
\texttt{I'd like to counter your offer by suggesting that I keep the two potatoes and give you the glue and one of the carrots in exchange for the berry. This way, I get to keep the potatoes, which I value highly, and you get the glue and one of the carrots, which you seemed to be interested in.} \\
            }
        }
    }
    & & \\ \\

    \theutterance \stepcounter{utterance}  
    & & \multicolumn{4}{p{0.6\linewidth}}{
        \cellcolor[rgb]{0.9,0.9,0.9}{
            \makecell[{{p{\linewidth}}}]{
                \texttt{\tiny{[P1$\langle$GM]}}
                \texttt{I think the glue and one of the carrots is a good starting point, but I'm not sure about giving up the two potatoes. I value the potatoes highly, and I'm not sure if I'm willing to part with them yet. } \\
\\ 
\texttt{I'd like to counter your offer by suggesting that I keep the two potatoes and give you the glue and one of the carrots in exchange for the berry. This way, I get to keep the potatoes, which I value highly, and you get the glue and one of the carrots, which you seemed to be interested in.} \\
            }
        }
    }
    & \\ \\

    \theutterance \stepcounter{utterance}  
    & \multicolumn{4}{p{0.6\linewidth}}{
        \cellcolor[rgb]{0.9,0.9,0.9}{
            \makecell[{{p{\linewidth}}}]{
                \texttt{\tiny{[P1$\rangle$GM]}}
                \texttt{I see your point about the potatoes. I'm willing to consider your counteroffer, but I'm not sure about giving up the berry. I value the berry more than the glue, so I'd like to propose a slight modification to your offer. How about I give you the glue and one of the carrots in exchange for the berry and one of the potatoes? This way, I get to keep one of the potatoes, which I value, and you get the glue and one of the carrots, as well as the berry.} \\
            }
        }
    }
    & & \\ \\

    \theutterance \stepcounter{utterance}  
    & & & \multicolumn{2}{p{0.3\linewidth}}{
        \cellcolor[rgb]{0.9,0.9,0.9}{
            \makecell[{{p{\linewidth}}}]{
                \texttt{\tiny{[GM$|$GM]}}
                \texttt{continue} \\
            }
        }
    }
    & & \\ \\

    \theutterance \stepcounter{utterance}  
    & & & \multicolumn{2}{p{0.3\linewidth}}{
        \cellcolor[rgb]{0.9,0.9,0.9}{
            \makecell[{{p{\linewidth}}}]{
                \texttt{\tiny{[GM$|$GM]}}
                \texttt{I see your point about the potatoes. I'm willing to consider your counteroffer, but I'm not sure about giving up the berry. I value the berry more than the glue, so I'd like to propose a slight modification to your offer. How about I give you the glue and one of the carrots in exchange for the berry and one of the potatoes? This way, I get to keep one of the potatoes, which I value, and you get the glue and one of the carrots, as well as the berry.} \\
            }
        }
    }
    & & \\ \\

    \theutterance \stepcounter{utterance}  
    & & \multicolumn{4}{p{0.6\linewidth}}{
        \cellcolor[rgb]{0.9,0.9,0.9}{
            \makecell[{{p{\linewidth}}}]{
                \texttt{\tiny{[GM$\rangle$P2]}}
                \texttt{I see your point about the potatoes. I'm willing to consider your counteroffer, but I'm not sure about giving up the berry. I value the berry more than the glue, so I'd like to propose a slight modification to your offer. How about I give you the glue and one of the carrots in exchange for the berry and one of the potatoes? This way, I get to keep one of the potatoes, which I value, and you get the glue and one of the carrots, as well as the berry.} \\
            }
        }
    }
    & \\ \\

    \theutterance \stepcounter{utterance}  
    & & & \multicolumn{4}{p{0.6\linewidth}}{
        \cellcolor[rgb]{0.9,0.9,0.9}{
            \makecell[{{p{\linewidth}}}]{
                \texttt{\tiny{[GM$\langle$P2]}}
                \texttt{I think your modification is a good compromise. I'm willing to accept your offer of giving me the glue and one of the carrots in exchange for the berry and one of the potatoes. This way, I get the items I value most, and you get to keep one of the potatoes, which is a good compromise for you. } \\
\\ 
\texttt{{[}Proposal: 1 glue, 1 carrot, 1 berry, 1 potato{]}} \\
            }
        }
    }
     \\ \\

    \theutterance \stepcounter{utterance}  
    & & & \multicolumn{2}{p{0.3\linewidth}}{
        \cellcolor[rgb]{0.9,0.9,0.9}{
            \makecell[{{p{\linewidth}}}]{
                \texttt{\tiny{[GM$|$GM]}}
                \texttt{proposal} \\
            }
        }
    }
    & & \\ \\

    \theutterance \stepcounter{utterance}  
    & & & \multicolumn{2}{p{0.3\linewidth}}{
        \cellcolor[rgb]{0.9,0.9,0.9}{
            \makecell[{{p{\linewidth}}}]{
                \texttt{\tiny{[GM$|$GM]}}
                [1, 1, 1, 1]
            }
        }
    }
    & & \\ \\

    \theutterance \stepcounter{utterance}  
    & & \multicolumn{4}{p{0.6\linewidth}}{
        \cellcolor[rgb]{0.9,0.9,0.9}{
            \makecell[{{p{\linewidth}}}]{
                \texttt{\tiny{[P1$\langle$GM]}}
                \texttt{The other player has submitted a secret proposal. It is now your turn to submit a proposal of your own.} \\
            }
        }
    }
    & \\ \\

    \theutterance \stepcounter{utterance}  
    & \multicolumn{4}{p{0.6\linewidth}}{
        \cellcolor[rgb]{0.9,0.9,0.9}{
            \makecell[{{p{\linewidth}}}]{
                \texttt{\tiny{[P1$\rangle$GM]}}
                \texttt{{[}Proposal: 1 berry, 2 carrots{]}} \\
            }
        }
    }
    & & \\ \\

    \theutterance \stepcounter{utterance}  
    & & & \multicolumn{2}{p{0.3\linewidth}}{
        \cellcolor[rgb]{0.9,0.9,0.9}{
            \makecell[{{p{\linewidth}}}]{
                \texttt{\tiny{[GM$|$GM]}}
                \texttt{proposal} \\
            }
        }
    }
    & & \\ \\

    \theutterance \stepcounter{utterance}  
    & & & \multicolumn{2}{p{0.3\linewidth}}{
        \cellcolor[rgb]{0.9,0.9,0.9}{
            \makecell[{{p{\linewidth}}}]{
                \texttt{\tiny{[GM$|$GM]}}
                [0, 0, 1, 2]
            }
        }
    }
    & & \\ \\

    \theutterance \stepcounter{utterance}  
    & & & \multicolumn{2}{p{0.3\linewidth}}{
        \cellcolor[rgb]{0.9,0.9,0.9}{
            \makecell[{{p{\linewidth}}}]{
                \texttt{\tiny{[GM$|$GM]}}
                [[0, 0, 1, 2], [1, 1, 1, 1]]
            }
        }
    }
    & & \\ \\

    \theutterance \stepcounter{utterance}  
    & \multicolumn{4}{p{0.6\linewidth}}{
        \cellcolor[rgb]{0.95,0.95,0.95}{
            \makecell[{{p{\linewidth}}}]{
                \texttt{\tiny{[P2$\rightarrow$P1]}}
                \texttt{Post{-}episode Reflexion (not sent to GM):} \\
\texttt{New plan: } \\
\\ 
\texttt{In my last attempt, I failed to achieve a mutually beneficial agreement with the other player. I made a secret proposal without confirming the agreement, which was against the rules. Moreover, my proposals were not complementary, as I gave up the berry, which I valued more than the glue. } \\
\\ 
\texttt{Next time, I will pay closer attention to the value function of the other player and try to make proposals that are complementary to theirs. I will also make sure to confirm the agreement with the other player before making a secret proposal. To achieve this, I will ask the other player to confirm their agreement with my proposal before I make a secret proposal of my own. For example, I could say "Are you agreeable to the proposal of giving you the glue and one of the carrots in exchange for the berry and one of the potatoes?" If they agree, I will make a secret proposal. If they do not agree, I will try to negotiate further to find a mutually beneficial agreement.} \\
            }
        }
    }
    & & \\ \\

    \theutterance \stepcounter{utterance}  
    & & & \multicolumn{4}{p{0.6\linewidth}}{
        \cellcolor[rgb]{0.95,0.95,0.95}{
            \makecell[{{p{\linewidth}}}]{
                \texttt{\tiny{[P1$\rightarrow$P2]}}
                \texttt{Post{-}episode Reflexion (not sent to GM):} \\
\texttt{New plan: } \\
\\ 
\texttt{In my last attempt, I made a mistake by not considering the value function of the other player when making a counteroffer. I also failed to think about the overall distribution of items and how it would affect the other player's proposal. Next time, I will make sure to consider the value function of the other player when making a counteroffer and think about the overall distribution of items.} \\
\\ 
\texttt{To achieve this, I will pay closer attention to the other player's value function and try to understand what they value most. I will also think about how my counteroffer will affect the overall distribution of items and try to find a compromise that benefits both players.} \\
\\ 
\texttt{For example, if the other player values the potatoes highly and I value the berry highly, I might consider making a counteroffer that involves trading the berry for one of the potatoes. This way, both players get to keep the items they value most, and the overall distribution of items is more balanced.} \\
\\ 
\texttt{I will also make sure to think about the other player's previous messages and try to understand their strategy. If they are willing to give up the berry, it might be because they value the glue and one of the carrots more. In this case, I might consider making a counteroffer that involves trading the berry for the glue and one of the carrots.} \\
            }
        }
    }
     \\ \\

\end{supertabular}
}

\end{document}
